\section{Justificación}

La formación práctica en ciberseguridad constituye una necesidad relevante para los estudiantes de Ingeniería Informática e Ingeniería de Sistemas de la Universidad Mayor de San Simón, debido a que las amenazas digitales actuales requieren no solo conocimiento teórico, sino también capacidad de análisis, toma de decisiones y respuesta oportuna ante incidentes. En la asignatura de Seguridad de Sistemas, el elevado número de estudiantes y el enfoque predominantemente conceptual dificultan el acompañamiento individualizado y la práctica guiada, generando una brecha entre lo aprendido en aula y su aplicación en situaciones cercanas al contexto profesional.

Si esta situación no se atiende, los estudiantes pueden finalizar su formación con dificultades para identificar riesgos digitales, priorizar medidas de protección y comprender las consecuencias de una respuesta inadecuada frente a incidentes como phishing, malware o ingeniería social. Esta limitación afecta el desarrollo de competencias prácticas en un área crítica para la protección de información, sistemas e infraestructuras tecnológicas, además de restringir las posibilidades de evaluación y retroalimentación inmediata dentro del proceso educativo.

En este contexto, el presente proyecto propone el desarrollo de un videojuego de simulación gamificada como herramienta tecnológica de apoyo al aprendizaje práctico de la ciberseguridad. A través de escenarios interactivos, retos progresivos, toma de decisiones y retroalimentación inmediata, el videojuego permitirá representar situaciones de amenaza en un entorno controlado, seguro y accesible, donde el estudiante pueda experimentar, equivocarse, corregir y fortalecer sus habilidades de respuesta sin comprometer sistemas reales.

Los principales beneficiarios serán los estudiantes que cursan contenidos relacionados con seguridad informática, ya que contarán con un recurso complementario para practicar conceptos y decisiones de defensa ante amenazas digitales. Asimismo, los docentes podrán disponer de una herramienta de apoyo pedagógico para dinamizar la enseñanza, observar el desempeño de los estudiantes y reforzar contenidos mediante actividades interactivas. De manera indirecta, la institución se beneficiará al incorporar una alternativa tecnológica orientada a mejorar la calidad de la formación práctica en ciberseguridad.

El proyecto es viable porque se desarrollará como un prototipo monojugador de alcance acotado, centrado en amenazas digitales fundamentales y en mecánicas de simulación educativa. Además, existen motores de desarrollo de videojuegos, recursos de diseño interactivo y herramientas tecnológicas accesibles que permiten implementar escenarios, retroalimentación y seguimiento de progreso dentro del tiempo y los recursos disponibles para Taller de Grado I.
